\documentclass[11pt]{article}

\usepackage[margin=1in]{geometry}
\usepackage{amsmath,amssymb,amsthm}
\usepackage{graphicx}
\usepackage[hidelinks]{hyperref}

\newtheorem{theorem}{Theorem}
\newtheorem{lemma}{Lemma}
\newtheorem{corollary}{Corollary}

\newcommand{\R}{\mathbb R}
\newcommand{\C}{\mathbb C}
\newcommand{\one}{\mathbf 1}
\newcommand{\spn}{\operatorname{span}}

\title{An \(L^1\) Obstruction to Naive Stone--von Neumann Uniqueness for Weyl Pairs}
\author{arXiv automatic math run \texttt{fa\_banach\_001}}
\date{June 28, 2026}

\begin{document}
\maketitle

\section*{Status}

This packet gives a candidate partial obstruction to Open Problem (4) of
van Neerven--Portal \cite{vNP2018}. It does not prove a Banach-space
Stone--von Neumann theorem. It shows that the most literal Hilbert-style
analogue, where irreducible Weyl pairs are expected to be similar to Hilbert
standard-pair models, is impossible in the full class of Banach spaces.

\section*{Source Question}

At the end of the paper the authors list four open problems:

\begin{center}
\includegraphics[width=\textwidth]{figures/open_problem_crop.png}
\end{center}

The target here is item (4): prove an analogue of the Stone--von Neumann
uniqueness theorem for Weyl pairs.

\section*{Definitions Used}

For the standard position/momentum pair on \(L^1(\R^d)\), write
\[
  (M_u f)(x)=e^{iu\cdot x}f(x),\qquad
  (T_v f)(x)=f(x+v),
  \qquad u,v,x\in\R^d.
\]
These are the bounded \(C_0\)-groups \(e^{iu\cdot Q}\) and \(e^{iv\cdot P}\)
from Example 3.2 of the source. A closed subspace \(Y\subseteq L^1(\R^d)\) is
called invariant for the standard Weyl pair if \(M_uY\subseteq Y\) and
\(T_vY\subseteq Y\) for every \(u,v\in\R^d\). The pair is topologically
irreducible if the only closed invariant subspaces are \(\{0\}\) and
\(L^1(\R^d)\).

\section*{Idea of the Proof}

The invariant-subspace part is a time-frequency cyclicity argument. If a
nonzero \(f\in L^1\) is fixed, then the closed span of all modulates and
translates \(M_uT_vf\) is all of \(L^1\). Indeed, any annihilator \(h\in
L^\infty\) would make the Fourier transform of \(x\mapsto f(x+v)h(x)\) vanish
for every translation \(v\), so \(f(x+v)h(x)=0\) almost everywhere for every
\(v\). A Fubini intersection argument then forces \(h=0\). This proves
irreducibility.

The obstruction to Hilbert-style uniqueness is then purely Banach-geometric:
\(L^1(\R^d)\) is non-reflexive, while any Hilbert direct sum of standard
\(L^2\)-models is reflexive.

\begin{lemma}[Cyclicity of one nonzero vector]
Let \(0\ne f\in L^1(\R^d)\). Then
\[
  \overline{\spn}\{M_uT_vf:\ u,v\in\R^d\}=L^1(\R^d).
\]
\end{lemma}

\begin{proof}
Let \(Z\) denote the closed span on the left. Suppose \(h\in L^\infty(\R^d)\)
annihilates \(Z\), in the usual duality
\[
  \langle g,h\rangle=\int_{\R^d}g(x)h(x)\,dx.
\]
Then for all \(u,v\in\R^d\),
\[
  0=\langle M_uT_vf,h\rangle
   =\int_{\R^d} e^{iu\cdot x} f(x+v)h(x)\,dx.
\]
For each fixed \(v\), the function \(x\mapsto f(x+v)h(x)\) belongs to
\(L^1(\R^d)\). By uniqueness of the Fourier transform on \(L^1(\R^d)\), it
follows that
\[
  f(x+v)h(x)=0
\]
for almost every \(x\), for every \(v\in\R^d\).

Assume, toward a contradiction, that \(h\ne0\). Choose \(\alpha,\beta>0\) and
a ball \(B_R\subset\R^d\) such that the sets
\[
  E=\{y:\ |f(y)|>\alpha\},\qquad
  F=\{x\in B_R:\ |h(x)|>\beta\}
\]
have positive finite measure. This is possible because \(f\ne0\) in \(L^1\)
and \(h\ne0\) in \(L^\infty\). By Fubini's theorem,
\[
  \int_{\R^d} |F\cap(E-v)|\,dv=|F|\,|E|>0.
\]
Hence for some \(v\), the set \(F\cap(E-v)\) has positive measure. On this set
\(|h(x)|>\beta\) and \(|f(x+v)|>\alpha\), contradicting
\(f(x+v)h(x)=0\) almost everywhere. Thus \(h=0\).

The annihilator of \(Z\) in \(L^\infty\) is zero. By Hahn--Banach, \(Z=L^1\).
\end{proof}

\begin{theorem}
For every \(d\ge1\), the standard position/momentum Weyl pair on
\(L^1(\R^d)\) is topologically irreducible.
\end{theorem}

\begin{proof}
Let \(Y\subseteq L^1(\R^d)\) be a nonzero closed subspace invariant under all
\(M_u\) and \(T_v\). Choose \(0\ne f\in Y\). Then \(M_uT_vf\in Y\) for all
\(u,v\in\R^d\), and by the lemma their closed linear span is all of
\(L^1(\R^d)\). Hence \(Y=L^1(\R^d)\).
\end{proof}

\begin{corollary}[Obstruction to a naive Hilbert-style analogue]
There is an irreducible Banach-space Weyl pair which is not similar to a direct
sum of Hilbert-space standard position/momentum pairs.
\end{corollary}

\begin{proof}
Take the standard pair on \(L^1(\R^d)\). By Example 3.2 of the source it is a
Weyl pair, and by the theorem it is topologically irreducible.

The Banach space \(L^1(\R^d)\) is non-reflexive. For instance, it contains a
closed isometric copy of \(\ell^1\), obtained from functions supported on
pairwise disjoint measurable sets of finite positive measure. Since \(\ell^1\)
is non-reflexive, \(L^1(\R^d)\) is non-reflexive as well. Any Hilbert direct
sum of standard \(L^2\)-models is a Hilbert space, hence reflexive. Therefore
there is no Banach-space isomorphism from \(L^1(\R^d)\) onto such a Hilbert
direct sum, and consequently no similarity intertwining the corresponding
Weyl pairs.
\end{proof}

\section*{Limitations}

This is only a partial obstruction. It does not rule out a Stone--von Neumann
analogue formulated inside a restricted Banach category, such as Hilbert
spaces, certain \(L^p\)-spaces with fixed \(p\), reflexive UMD spaces, Banach
lattices with extra structure, or non-commutative \(L^p\)-spaces. It also does
not address the vector-valued twisted-convolution problems (1)--(3).

\section*{Novelty Check}

The local indexes were searched for arXiv:1806.00980, Weyl calculus, Weyl
pairs, group generators, canonical commutation relations, Stone--von Neumann,
twisted convolution, UMD, and functional-calculus keywords. No exact prior
packet for this open-problem item was present.

External searches on June 28, 2026 for combinations of the source title,
Stone--von Neumann, Weyl pairs, Banach spaces, twisted convolution, and UMD did
not locate a later solution to the source open problems. Since this packet is
a partial obstruction rather than a full solution, the novelty claim is modest:
it records a simple boundary condition that any positive analogue must respect.

\section*{Human Review Recommendation}

Verify that topological irreducibility is an acceptable Banach-space analogue
of Hilbert irreducibility for the intended discussion. If the intended open
problem already restricts to Hilbertizable or reflexive settings, this packet
should be read only as a scope warning. Under the literal wording, it is a
valid obstruction to a naive Hilbert-model uniqueness theorem.

\begin{thebibliography}{9}

\bibitem{vNP2018}
J. van Neerven and P. Portal,
\emph{The Weyl calculus for group generators satisfying the canonical
commutation relations},
arXiv:1806.00980.

\end{thebibliography}

\end{document}
